% Paper 24 §V drop-in: Spinor-Lichnerowicz corollary
% Insert as a new subsection after §V.B (Two-fermion entanglement rigidity)
% Track T9, Tier 3 sprint, 2026-04-15

\subsection{Spinor-Lichnerowicz corollary: the operator-order discriminant}
\label{sec:spinor-lichnerowicz}

The HO rigidity theorem of Section~\ref{sec:ho-rigidity-theorem}
and the Coulomb/HO asymmetry of Section~\ref{sec:coulomb-ho-asymmetry}
established that calibration~$\pi$ (even-zeta content) enters from
second-order Riemannian operators on scalar bundles: the HO via
the Bargmann--Segal Euler operator on $H^2(S^5)$, and the Coulomb
system via the Fock projection to $-\Delta_{\mathrm{LB}}$ on~$S^3$.
The first-order/second-order distinction was identified as the
structural origin of the $\pi$-free property of the Bargmann graph
(first-order, complex-analytic, linear projection) versus the
calibration~$\pi$ of the Fock Dirichlet series (second-order,
Riemannian, nonlinear projection).

The Dirac-on-$S^3$ Tier~3 sprint (Track~T9) extends this result to
spinor bundles.  The squared Dirac operator on unit~$S^3$ satisfies
the Lichnerowicz identity $D^2 = \nabla^*\nabla + R/4$
\cite{lichnerowicz1963}, where $R = 6$.  Its spectral zeta evaluates
at every integer~$s$ to a two-term polynomial in~$\pi^2$ with
rational coefficients (Paper~18, Eq.~\ref{eq:zeta_d2}), with
\emph{no odd-zeta content} at any~$s$.

The mechanism is algebraic: squaring maps the first-order Dirac
sum $\sum g_n \cdot (\text{odd})^{-s}$ (which produces
$\zeta_R(\text{odd})$ at odd~$s$, as in Tier~1 Track~D3) to the
second-order sum $\sum g_n \cdot (\text{odd})^{-2s}$, which is always
$\mathbb{Q} \cdot \pi^{2s}$ by the Bernoulli formula.  The
odd-zeta content of~$D$ is structurally lost upon squaring.

\begin{corollary}[Operator-order transcendental discriminant]
The operator-order distinction (first-order vs.\ second-order) is
bundle-independent on round~$S^3$.  All second-order operators---scalar
Laplace--Beltrami, spinor connection Laplacian $\nabla^*\nabla$,
squared Dirac~$D^2$---produce exclusively $\pi^{\mathrm{even}}$
spectral-zeta content.  All first-order operators (Dirac~$D$) can
produce odd-zeta content ($\zeta_R(3), \zeta_R(5), \ldots$) via
half-integer eigenvalue sums.  The bundle type (scalar vs.\ spinor)
modulates rational coefficients and the eigenvalue/charge lattice
(integer vs.\ half-integer) but does not change the transcendental
class.
\end{corollary}

This corollary completes the universal/Coulomb-specific partition
stated in the abstract: the Bargmann graph is $\pi$-free because
it encodes a first-order operator (the Euler operator); calibration
$\pi$ is second-order; and this pattern persists on spinor bundles.
Cross-reference: Paper~18 \S IV (taxonomy closure via
Eq.~\ref{eq:zeta_d2}); Tier~1 Track~D3 ($\zeta(3)$ from
first-order Dirac); T9 debug memo
(\texttt{debug/dirac\_t9\_memo.md}).
